\documentclass[varwidth=15cm,border=10pt]{standalone}
\usepackage{amsmath,xcolor,array}
\usepackage{fontspec,unicode-math}
\setmainfont{DejaVu Sans}[ItalicFont={DejaVu Sans},ItalicFeatures={FakeSlant=0.2},BoldItalicFont={DejaVu Sans Bold},BoldItalicFeatures={FakeSlant=0.2}]
\setmathfont{texgyredejavu-math.otf}
\setmathfont{DejaVu Sans}[range={up,bfup}]
\setmathfont{DejaVu Sans}[range={it,bfit},FakeSlant=0.2]
\setmathfont{texgyredejavu-math.otf}[range={"2202,"1D715,"03B5,"03F5,"2112,"2131}]
\pagecolor[HTML]{1E1E2E}\color[HTML]{F4F4F4}
\definecolor{accent}{HTML}{89B4FA}
\definecolor{warn}{HTML}{F38BA8}
\definecolor{good}{HTML}{A6E3A1}
\definecolor{ref}{HTML}{9399B2}
\newcommand{\heading}[1]{\par\medskip{\large\bfseries\color{accent}#1}\par\smallskip}
\newcommand{\src}[1]{{\color{ref}\small[#1]}}
\newcommand{\dolabel}{{\bfseries DO}\ \ }
\newcommand{\why}{{\bfseries\color{accent}WHY}\ \ }
\newcommand{\chk}{{\bfseries\color{good}CHECK}\ \ }
\newcommand{\rep}{{\bfseries\color{warn}WRITE}\ \ }
\setlength{\parindent}{0pt}
\raggedright
\begin{document}

{\Large\bfseries Solving A3 Question 2 yourself — the procedure}

\heading{Before you start}
L06 slide 24 lists five checks and then says: {\itshape ``This is the workflow
Assignment 3 asks you to carry out.''} Question 2 \emph{is} that list. Slide 33
adds the other half: {\itshape ``A different structure is a different boundary
condition, not different code.''}

{\bfseries The one idea.} You are not writing a truss program. You are changing
one block of the lecture's program, then reading the answer with a fixed
checklist. The marks are in the reading, not in the running.

\heading{Step 1 — Run the lecture demo unchanged, first}
\dolabel \texttt{l06\_demo\_truss\_sizeOpt.m}, as supplied, on its own
$7\times3$ grid over $6\times2$~m.\par
\why If your edited version misbehaves later you need to know whether you broke
it or whether it was never working here. Thirty seconds now.

\heading{Step 2 — Change the ONE block: the boundary conditions}
\dolabel Replace the whole-left-edge support with two pinned points, and move the
load:
\[
  \begin{aligned}
    &\text{pinned, both directions, at } (0,\ L_y/4) \text{ and } (0,\ 3L_y/4)\\
    &P = 10^{5}\,\text{N, downward, at } (L_x,\ L_y/2)
  \end{aligned}
\]
Specify the heights \emph{physically} and snap each to the nearest left-edge
node.\par
\why The assignment says to, so that the same code keeps working when the grid
changes in 2.2. At $h_y = 1$~m and $0.5$~m both snaps are exact.\par
{\color{warn}Slides vs assignment --- do not copy the slide.} Slide 33 describes
this task as supports at $L_y/3$ and $2L_y/3$ on a $6\times3$~m domain. The
assignment PDF says $L_y/4$ and $3L_y/4$ on $6\times4$~m. {\bfseries The PDF is
what you are marked against.}

\heading{Step 3 — Enter the baseline data exactly}
\small
\begin{tabular}{@{}ll@{}}
domain & $L_x = 6$~m, $L_y = 4$~m\\
grid & $n_x = 7$, $n_y = 5$, \texttt{'full'} ($h_x = h_y = 1$~m)\\
budget & $V = 0.2$~m$^3$\\
material & $E = 2.1\times10^{11}$~Pa\\
bounds & $A_{\min} = 10^{-7}$, $A_{\max} = 10^{-2}$~m$^2$\\
\end{tabular}
\normalsize
\par\smallskip
\why $A_{\min}$ is $10^{-7}$ and not $0$ so that a bar whose area vanishes still
contributes something: if every bar at a node went to zero, $K$ would be
singular.

\heading{Step 4 — Solve the non-dimensionalised problem, never the physical one}
\dolabel Three steps, in order \src{slides 12--14}: parameterise the area,
non-dimensionalise the objective, normalise the budget.
\[
  A_e = A_{\min} + (A_{\max}-A_{\min})x_e,
  \qquad
  \hat f(x) = \frac{f(a(x))}{f_0},
  \qquad
  Ax\le b,\quad 0\le x_e\le 1 .
\]
\why Not cosmetic. Slide 15 reports the same problem solved both ways:
non-dimensionalised exits with {\bfseries flag 1}, the original with areas
directly exits with {\bfseries flag 2}, step size gone to zero. {\itshape ``It
decides whether the run converges.''} Solve in $x$, report in $a$.

\heading{Step 5 — Item 1: count the fixed DOFs before you solve}
\dolabel \textbf{4 of 70}: two pinned points $\times$ two directions.\par
\why The question asks what the number \emph{should} be, so justify it: only two
\emph{points} are pinned, not an edge. Three is the minimum that removes two
translations and one rotation, so four is just-stiff with one redundancy. The
lecture's whole-edge support would give 10 at $n_y=5$ --- if you print 10, you
did not change the block.

\heading{Step 6 — Items 2 and 3: read three numbers, in this order}
\dolabel \texttt{exitflag}, then \texttt{output.iterations}, then
\texttt{output.firstorderopt} \src{slide 25}.\par
\why {\itshape ``\texttt{exitflag} --- the only one that says why it stopped.
Anything but 1 and the rest is premature.''} Iterations close to
\texttt{MaxIterations} means it ran out of patience, not that it converged.
\texttt{firstorderopt} is the KKT residual actually reached, to be quoted against
\texttt{OptimalityTolerance}.\par
\rep ``Flag 1: first-order optimality and feasibility met to tolerance'' --- and
the two numbers that back it.

\heading{Step 7 — Items 4 and 5: the objective and the material}
\dolabel $f_0$, $f^\star$, $f^\star/f_0$; then the volume \emph{both} ways ---
physical $\sum_e L_eA_e$ in m$^3$, and the normalised constraint value
$\left(\sum_e L_eA_e\right)/V$.\par
\chk The normalised one must sit at $1$ to several decimals. If it does not, the
budget is not active and something is wrong --- see the next step.

\heading{Step 8 — Item 6: the multiplier is a PREDICTION, not a surprise}
\dolabel Report \texttt{lambda.ineqlin} and state that it is positive.\par
\why Slide 18 derives it before any run:
\[
  \frac{\partial f}{\partial A_e} = -u_e^\top K_e^{\text{unit}}u_e \le 0 ,
\]
twice a strain energy with a minus sign. Thickening \emph{any} bar can only
reduce compliance, and nothing in the objective stops the optimiser thickening
everything --- so the budget must be active and $\lambda>0$. The slide spells out
the consequence: {\itshape ``in Assignment 3 the volume should print as fully
spent and $\lambda>0$ in every run --- if it does not, something is wrong with
your constraint.''}\par
{\color{good}Free extra check:} for a fixed layout, compliance is exactly
proportional to $1/V$, so $\lambda_{\text{ineqlin}}$ should come out equal to
$f^\star/f_0$. Agreement tests the constraint scaling, the objective
normalisation and the multiplier at once.

\heading{Step 9 — Item 7: how many bars survive}
\dolabel Count bars above 1\% of the largest area, and say what fraction of the
volume they hold.\par
\why This is Check 3, ``where did the material go?'' \src{slide 26}. The point to
make is slide 31's: {\bfseries a size optimisation problem, with the nodes fixed
in advance, has chosen a topology} --- by driving most bars onto $A_{\min}$.

\heading{Step 10 — Item 8: three figures, and how to read them}
\dolabel The optimised structure, the convergence history, the area histogram.\par
\why Structure is Check 4, {\itshape ``observe the optimised layout and name the
load path''} \src{slide 27}. Expect a tapered lens --- a Michell cantilever
\src{slide 28}.\par
{\color{warn}The convergence trap:} the staircase of flat stretches and sudden
drops is the interior-point barrier parameter $\mu$ being cut \src{slides
22--23}. A plateau is one barrier subproblem converging --- {\bfseries it is not
convergence}. That is why the exit flag is read before the curve.\par
The histogram should be bimodal: a crowd at $A_{\min}$, a few thick bars, little
between. The optimiser discards bars rather than thinning them.

\heading{Step 11 — 2.2: run A and B through the SAME solver}
\dolabel Baseline $(7,5)$ \texttt{'full'}; A $(7,5)$ \texttt{'neighbour'}; B
$(13,9)$ \texttt{'full'}. Everything not named stays at baseline.\par
\dolabel \textbf{Item 1:} the three structures side by side, {\bfseries same
bar-width scale} and {\bfseries axes fixed to the full $L_x\times L_y$ domain in
every panel}.\par
\why The assignment gives the reason: otherwise a given thickness means a
different area in each picture, and a design that does not reach the edge gets
stretched to fill its own panel. Both halves are marked.

\heading{Step 12 — Item 2: fill the table, then say what it does NOT mean}
{\color{warn}The trap:} $f_0$ differs between runs by design --- the same budget
spread over different total bar lengths. So $f^\star/f_0$ and $\lambda$ are
\emph{within-run} quantities. {\bfseries Only $f^\star$ compares across runs.}

\heading{Step 13 — Item 3: connectivity, and why it MUST be that one}
\dolabel \texttt{'full'} beats \texttt{'neighbour'}.\par
\why Set inclusion, not observation: every \texttt{'neighbour'} design is also a
\texttt{'full'} design with the long diagonals left at $A_{\min}$, so
\texttt{'neighbour'}'s feasible set is a \emph{subset}, and a minimum over a
subset cannot be lower. Guaranteed before either run started.\par
\rep ``What do you pay for it?'' is part of the question: bars and seconds.

\heading{Step 14 — Item 4: refinement, and the trap the lecture sets}
\dolabel B beats the baseline.\par
\why Only because the grids \emph{nest}. Halving both spacings keeps every
coarse node. A coarse bar is not itself a candidate in the fine structure --- the
\texttt{'full'} rule drops any bar with a third node on it --- but two collinear
fine bars of length $L/2$ and area $A$ in series have stiffness $EA/L$ and volume
$LA$, exactly the coarse bar. So the coarse feasible set embeds in the fine one
and $f^\star$ must fall.\par
{\color{warn}Do not write ``refinement improves $f^\star$''.} The lecture's own
table \src{slide 29} has $11\times5 \to 13\times5$ \texttt{'neighbour'} going
$2.220\times10^{2} \to 2.302\times10^{2}$ --- {\bfseries worse}. Those grids do
not nest. Copy the lecture's conclusion and you answer backwards.\par
\rep Cost in bars and in run time, and whether you would pay it. The question
asks all three.

\heading{Appendix — the five checks, and what to cite}
\src{slide 24} 1. Did it converge --- the exit flag. 2. Read the multiplier.
3. How is the material distributed --- the histogram. 4. Observe the layout and
name the load path. 5. Effect of ground structure and refinement.
\par\smallskip
\begin{tabular}{@{}p{5.0cm}p{9.2cm}@{}}
Rao \textit{\S12.4}, p.~622 & with \S12.4.1--12.4.3 (pp.~623--625): truss sizing.\\[3pt]
Rao \textit{\S15.4}, p.~705 & design sensitivities, from Lecture 5 --- reread ``as the tool that will be used''.\\[3pt]
L06 slides 1--33 & formulation 11--15, sensitivity 16--18, reading the solver 19--23, the five checks 24--31.\\
\end{tabular}

\end{document}
